% ============================================================
\section{Code-Design: SOLID \& Patterns (Abschnitt VIII)}
% ============================================================
\emph{Software Engineering} ist die systematische Entwicklung, Planung, Implementierung, Testung und Wartung von Software. Gutes Design reduziert \textbf{Kopplung} und f\"ordert \textbf{klare Verantwortlichkeiten}.

\subsection{SOLID-Prinzipien}
F\"unf Designprinzipien f\"ur skalierbare, erweiterbare und flexible (OO-)Software.

\subsubsection*{S -- Single Responsibility (SRP)}
Eine Klasse sollte \textbf{genau einen Grund zur \"Anderung} haben. Mehrere Aufgaben in einer Klasse $\Rightarrow$ \"Anderungen beeinflussen sich gegenseitig, schlechtere Test- und Wiederverwendbarkeit.
\begin{lstlisting}
// Verletzung: drei Verantwortlichkeiten          // Korrektur: aufteilen
class Report {                                     class Report       { void generate(); };
    void generate();                               class ReportSaver  { void saveToFile(const Report&); };
    void saveToFile();                             class ReportPrinter{ void print(const Report&); };
    void print();
};
\end{lstlisting}
\begin{merke}
,,Eine Aufgabe'' hei\ss t \textbf{nicht} ,,eine Methode'' -- eine Klasse darf viele Methoden haben, solange sie zur gleichen Verantwortung geh\"oren.
\end{merke}

\subsubsection*{O -- Open/Closed (OCP)}
Klassen sollen \textbf{offen f\"ur Erweiterung, aber geschlossen f\"ur \"Anderung} sein: neuen Code hinzuf\"ugen statt funktionierenden Code umschreiben (vermeidet Bugs/Regressionen). Umsetzung \"uber Interfaces, abstrakte Klassen, virtuelle Methoden, Polymorphie.
\begin{lstlisting}
// statt if (shape=="circle") ... else if ...   ->   Vererbung + override
class Shape  { public: virtual double area() const = 0; virtual ~Shape() = default; };
class Circle : public Shape { public: double area() const override { return 10.0; } };
\end{lstlisting}

\subsubsection*{L -- Liskov Substitution (LSP)}
Objekte einer abgeleiteten Klasse m\"ussen sich \textbf{wie} die der Basisklasse verhalten -- sie m\"ussen \"uberall austauschbar sein. Drei Regeln:
\begin{enumerate}
  \item Vorbedingungen d\"urfen \textbf{nicht versch\"arft} werden (Derived verlangt nicht mehr als Base).
  \item Nachbedingungen d\"urfen \textbf{nicht abgeschw\"acht} werden (Derived garantiert mind. so viel wie Base).
  \item Invarianten m\"ussen \textbf{erhalten} bleiben (global g\"ultige Zust\"ande nicht verletzen).
\end{enumerate}

\subsubsection*{I -- Interface Segregation (ISP)}
Klassen sollen nicht von Methoden abh\"angen, die sie nicht brauchen. Lieber \textbf{mehrere kleine, spezialisierte Interfaces} als ein gro\ss es. (ISP ist das SRP f\"ur Interfaces.)
\begin{lstlisting}
// Verletzung: BasicPrinter muss scan() implementieren, wirft aber nur Exception
class Machine { public: virtual void print() = 0; virtual void scan() = 0; };
// Besser: getrennte Interfaces Printer und Scanner
\end{lstlisting}

\subsubsection*{D -- Dependency Inversion (DIP)}
High-Level-Module sollen nicht von konkreten Implementierungen abh\"angen, sondern von \textbf{Abstraktionen} (Interfaces). Entkoppelt Fachlogik von technischen Details (DB, API, Framework).
\begin{lstlisting}
class Database { public: virtual void save_order() = 0; };       // Abstraktion
class MySQLDatabase : public Database { public: void save_order() override {} };
class OrderService {
    Database& db;                                                 // kennt nur die Abstraktion
public:
    OrderService(Database& db) : db(db) {}
};
\end{lstlisting}

\subsection{Design Patterns}
\begin{definition}[: Design Pattern]
Bew\"ahrte, wiederverwendbare L\"osungs\emph{baupl\"ane} f\"ur wiederkehrende Design-Probleme -- kein konkreter Code, sondern Beschreibung von Problem, L\"osungsansatz und beteiligten Komponenten. Standardwerk: \emph{Gang of Four} (Gamma et al., 1994).
\end{definition}
\begin{center}\small
\begin{tabular}{@{}lll@{}}
\toprule
\textbf{Kategorie} & \textbf{Zweck} & \textbf{Beispiele} \\
\midrule
Creational (Erzeugung) & Objekterstellung regeln/entkoppeln & Factory, Singleton, Builder \\
Structural (Struktur) & Klassenaufbau vereinfachen & Adapter, Decorator, Proxy \\
Behavioral (Verhalten) & Kommunikation/Verantwortung & Observer, Strategy, Command, State \\
\bottomrule
\end{tabular}
\end{center}

\subsubsection*{Factory Method (Creational)}
Objekte werden nicht \"uberall direkt mit \code{new} erzeugt, sondern zentral \"uber eine \emph{Factory} -- gekapselt, austauschbar, erweiterbar, testbar.
\begin{lstlisting}
class SensorFactory {
public:
    static std::unique_ptr<Sensor> create_sensor(const std::string& type) {
        if (type == "temperature") return std::make_unique<TemperatureSensor>();
        if (type == "pressure")    return std::make_unique<PressureSensor>();
        return nullptr;
    }
};
auto sensor = SensorFactory::create_sensor("temperature");
sensor->read();
\end{lstlisting}
Vorteil: \code{main} kennt nur noch die Factory + die abstrakte Basisklasse \code{Sensor} (DIP); neue Sensoren erfordern nur eine \"Anderung der Factory (OCP).

\subsubsection*{Observer (Behavioral)}
Ein \emph{Subject} informiert automatisch registrierte \emph{Observer}, wenn sich sein Zustand \"andert. Statt dass z.\,B. ein Sensor \code{display.update()}, \code{logger.update()}, \dots\ direkt aufruft, registrieren sich die Observer \"uber ein gemeinsames Interface $\Rightarrow$ lose Kopplung, leicht erweiterbar (OCP).

\subsubsection*{Adapter (Structural)}
\"Ubersetzt zwischen zwei inkompatiblen Schnittstellen: Eine fremde Klasse (z.\,B. \code{LegacyTemperatureDevice}) soll genutzt werden, ohne sie zu \"andern. Der Adapter sieht nach au\ss en wie das erwartete \code{Sensor}-Interface aus, nutzt intern aber die Fremdklasse. Ideal zum Integrieren von Legacy-/Herstellerbibliotheken.

\begin{tipp}
SOLID + Patterns zielen auf dasselbe: \textbf{geringe Kopplung, hohe Kohäsion}. Patterns schaffen au\ss erdem eine gemeinsame Sprache -- ,,das ist ein Strategy Pattern'' ersetzt lange Erkl\"arungen.
\end{tipp}
\clearpage
